\documentclass[11pt]{article}
\usepackage[T1]{fontenc}
\usepackage{lmodern}
\usepackage[margin=1in]{geometry}
\usepackage{amsmath,amssymb,amsthm,mathtools}
\usepackage{microtype,booktabs,array}
\usepackage[dvipsnames]{xcolor}
\usepackage[colorlinks=true,linkcolor=MidnightBlue,urlcolor=MidnightBlue,citecolor=MidnightBlue]{hyperref}
\usepackage{fancyhdr}
\pagestyle{fancy}\fancyhf{}
\fancyhead[L]{\small Depth-one Artin--Schreier defect towers}
\fancyhead[R]{\small Dev D. Goyal}
\fancyfoot[C]{\thepage}
\setlength{\headheight}{14pt}
\setlength{\parskip}{4pt}\setlength{\parindent}{0pt}
\allowdisplaybreaks
\newtheorem{theorem}{Theorem}
\newtheorem{lemma}{Lemma}[section]
\newtheorem{proposition}[lemma]{Proposition}
\newtheorem{corollary}{Corollary}
\newcommand{\code}[1]{\nolinkurl{#1}}
\newcommand{\Q}{\mathbb Q}
\newcommand{\R}{\mathbb R}
\newcommand{\F}{\mathbb F}
\newcommand{\supp}{\operatorname{supp}}
\newcommand{\cl}{\operatorname{cl}}
\newcommand{\trn}{\operatorname{trn}}
\newcommand{\depth}{\operatorname{depth}}
\newcommand{\Gal}{\operatorname{Gal}}
\hypersetup{pdftitle={Depth-one Artin-Schreier defect towers are composita},
pdfauthor={Dev D. Goyal},
pdfsubject={Fully autonomous agentic research},
pdfkeywords={Artin-Schreier extensions, defect, depth, Okutsu sequences, Hahn series, Krasner lemma, Lean 4}}
\newcommand{\replaystatus}{Fresh replays on 30 and 31 October 2026 compiled both Lean projects one module at a time into new directories, against Mathlib checkouts at the pinned revisions: the 20 modules of the Lean~4.19 project and the 100 modules of the Lean~4.35 project, 54 of them vendored, each with a regenerated audit. All 435 and 2127 declarations defined in them, including generated and private ones, depend only on \code{propext}, \code{Classical.choice} and \code{Quot.sound}. Both implications are formalized, with depth one defined by Mac Lane--Vaquié chains in the second project; Section~\ref{sec:lean} states the formal scope.}
\title{\textbf{Depth-one Artin--Schreier defect towers are composita}}
\author{Dev D. Goyal}
\date{30 October 2026}

\begin{document}
\maketitle

\begin{abstract}
Let $k$ be an algebraically closed field of characteristic $p>0$ and let $K$ be a perfect henselian subfield of the Hahn field $k((y^\Q))$ that contains $k(y)$ and whose elements have no finite accumulation points in their supports. Nart and Novacoski proved that every finite compositum of Artin--Schreier extensions of $K$ has depth one, and conjectured the converse. We prove the converse: a finite Artin--Schreier tower $M/K$ of depth one is a compositum of Artin--Schreier extensions of $K$. More precisely, for every $\theta\in M\setminus K$ the truncation of $\theta$ at the first finite accumulation point of its support generates a nontrivial elementary abelian subextension of $K(\theta)/K$, and depth one forces this subextension to be all of $M$. The extraction step adapts an Ore-relation argument of Kedlaya. Both implications are formalized in Lean~4 for actual Artin--Schreier towers over a base whose valuation ring is henselian, with depth one defined as the least length of a complete Mac Lane--Vaquié chain over all primitive generators.
\end{abstract}

\paragraph{Verification status.} \replaystatus

\section{The theorem}\label{sec:theorem}

Let $k$ be an algebraically closed field of characteristic $p>0$ and let $H=k((y^\Q))$ be the field of Hahn series with well-ordered rational support. Write $v$ for the valuation by least exponent. Suppose
\[
k(y)\subseteq K\subseteq H,
\]
that $K$ is perfect and henselian for $v$, and that every $a\in K$ has only finitely many support exponents below every real bound; for Hahn series this says that the support has no finite accumulation point. These are the standing assumptions of Nart and Novacoski~\cite[\S3]{NN}.

An \emph{Artin--Schreier tower} is a chain $K=K_0\subset K_1\subset\dots\subset K_r=M$ in which each step is generated by a root of $X^p-X-a_i$ with $a_i\in K_{i-1}$. Redundant steps can be removed, so we assume every step has degree $p$ and $r\ge1$. Such towers can be realized in $H$ by the explicit root formula of Lemma~\ref{lem:roots}, and by henselianity the valuation does not depend on the realization.

For an algebraic element $\theta$ write $\deg_K\theta=[K(\theta):K]$. Nart and Novacoski define the depth of a primitive generator through Mac Lane--Vaquié chains and characterize it through Okutsu sequences~\cite[Definition~1.2, Theorem~1.3]{NN}. For a nontrivial finite immediate extension, a primitive generator $\theta$ has depth one exactly when
\begin{equation}\label{eq:depthone}
\text{for every algebraic } b \text{ with } \deg_K b<\deg_K\theta \text{ there is } a\in K \text{ with } v(\theta-b)\le v(\theta-a).
\end{equation}
The depth of $M/K$ is the minimum depth over all primitive generators. As in~\cite[\S3]{NN}, distances are measured in the algebraic closure $\overline K$ of $K$ inside $H$; the argument below only uses approximations lying in $M$ itself.

\begin{theorem}\label{thm:main}
Under these assumptions, let $M/K$ be a nontrivial finite Artin--Schreier tower. If $\depth(M/K)=1$, then $M$ is a compositum of Artin--Schreier extensions of $K$.
\end{theorem}

Nart and Novacoski proved the reverse implication~\cite[Theorem~3.15]{NN} and conjectured Theorem~\ref{thm:main}~\cite[Conjecture~4.1]{NN}. Together they give:

\begin{corollary}
Under the standing assumptions, a nontrivial finite Artin--Schreier tower $M/K$ has depth one if and only if $M$ is a compositum of Artin--Schreier extensions of $K$.
\end{corollary}

The proof gives a stronger statement, Proposition~\ref{prop:block}: for every $\theta\in M\setminus K$, the first accumulation block of $\theta$ generates a nontrivial elementary abelian Galois subextension of $K(\theta)/K$.

\section{Monomials and supports of the base}\label{sec:base}

Put $\Gamma=vK\subseteq\Q$. Since $K$ is perfect, $\Gamma$ is $p$-divisible.

\begin{lemma}\label{lem:monomial}
For every $q\in\Gamma$, the monomial $y^q$ lies in $K$.
\end{lemma}

\begin{proof}
Choose $a\in K$ with $v(a)=q$ and leading coefficient $\lambda\in k^\times$, and write $q=m/(p^rn)$ with $p\nmid n$. Perfection gives $y^{m/p^r}\in K$, so $u=a^n/y^{m/p^r}\in K$ is a unit with residue $\lambda^n$. Hensel's lemma applied to $X^n-u$ at the simple root $\lambda$ of its reduction gives $b\in K$ with $b^n=u$ and residue $\lambda$. In $H$, $a/y^q$ is another $n$th root of $u$ with residue $\lambda$; their quotient is an $n$th root of unity, these all lie in $k$, and the only one with residue $1$ is $1$. Hence $b=a/y^q$ and $y^q=a/b\in K$.
\end{proof}

\begin{lemma}\label{lem:basesupport}
Every exponent in the support of an element of $K$ lies in $\Gamma$, and every finite truncation of an element of $K$ lies in $K$.
\end{lemma}

\begin{proof}
If an exponent outside $\Gamma$ occurs, let $q$ be the least one. The earlier terms form a finite sum, since $K$ has no finite accumulation points, and this sum lies in $K$ by Lemma~\ref{lem:monomial}. Subtracting it leaves an element of $K$ of valuation $q\notin\Gamma$, which is impossible. Finite truncations are then finite $k$-linear combinations of monomials in $K$.
\end{proof}

\section{Accumulation points in an Artin--Schreier tower}\label{sec:support}

Say that a Hahn series $s$ has property $C_\Gamma$ if the closure in $\R$ of its support is contained in $\Gamma$. The closure must be taken in $\R$: closure in $\Q$ would miss irrational accumulation points.

\begin{lemma}\label{lem:closure}
Property $C_\Gamma$ is preserved by sums and products.
\end{lemma}

\begin{proof}
The support of a sum lies in the union of the supports, and the closure of a finite union is the union of the closures. The support of a product $st$ lies in $S+T$ with $S=\supp s$ and $T=\supp t$, both bounded below. If $s_n+t_n$ converges with $s_n\in S$, $t_n\in T$, both coordinate sequences are bounded, and a subsequence gives limits in $\cl S$ and $\cl T$. Hence $\cl(S+T)\subseteq\cl S+\cl T\subseteq\Gamma$.
\end{proof}

\begin{lemma}\label{lem:roots}
If $a$ has property $C_\Gamma$, then every root in $H$ of $X^p-X-a$ has property $C_\Gamma$.
\end{lemma}

\begin{proof}
Write $a=a^-+a_0+a^+$ by negative, zero and positive exponents, and choose $c\in k$ with $c^p-c=a_0$. One root is
\[
\beta=\sum_{n\ge1}(a^-)^{1/p^n}+c-\sum_{n\ge0}(a^+)^{p^n}.
\]
Both sums are Hahn-summable. The supports of the negative terms lie in $[v(a^-)/p^n,0)$, so away from $0$ only finitely many $n$ contribute; if $a^+\ne0$ has least exponent $\varepsilon>0$, only finitely many $n$ with $p^n\varepsilon$ below a given bound contribute to the positive sum. Hence the union of the supports is well ordered and each exponent receives finitely many contributions, and applying Frobenius termwise gives $\beta^p-\beta=a$. The closure of the negative part lies in $\{0\}\cup\bigcup_{n\ge1}p^{-n}\cl(\supp a^-)$, and the positive part is locally finite on bounded intervals with closure in $\bigcup_np^n\cl(\supp a^+)$. Since $\Gamma$ is $p$-divisible and closed under multiplication by $p$, all these points lie in $\Gamma$. The other roots differ from $\beta$ by elements of $\F_p$.
\end{proof}

\begin{proposition}\label{prop:tower}
Every element of $M$ has property $C_\Gamma$. In particular its support and all its finite accumulation points lie in $\Gamma$, and $M/K$ is immediate.
\end{proposition}

\begin{proof}
Elements of $K$ have property $C_\Gamma$ by Lemma~\ref{lem:basesupport} and the absence of finite accumulation points. If $K_i=K_{i-1}(\beta)$, Lemma~\ref{lem:roots} gives $C_\Gamma$ for $\beta$, every element of $K_i$ is a polynomial in $\beta$ of degree less than $p$ over $K_{i-1}$, and Lemma~\ref{lem:closure} completes the induction. Consequently $vM\subseteq\Gamma=vK$, and the residue field of $M$ is $k$ because $k\subseteq K\subseteq M\subseteq H$ and $H$ has residue field $k$.
\end{proof}

The tower hypothesis is used here. In characteristic $2$ the series $s=\sum_{n\ge1}y^{(1-4^{-n})/3}$ satisfies $s^4+ys+y=0$ and has support in $\mathbb Z[1/2]$ but the accumulation point $1/3$, so support in $\Gamma$ alone does not control accumulation points.

\section{The first accumulation block}\label{sec:block}

Fix $\theta\in M\setminus K$ and put $n=\deg_K\theta>1$.

\begin{lemma}\label{lem:accumulation}
The support of $\theta$ has a finite accumulation point. If $\delta$ is the least one, then $\delta\in\Gamma$, the truncation
\[
\alpha=\trn_\delta(\theta)=\sum_{q<\delta}\theta_q\,y^q
\]
has infinite support, and every proper lower truncation of $\alpha$ is finite.
\end{lemma}

\begin{proof}
If the support had no finite accumulation point, the finite truncations of $\theta$ would lie in $K$ by Lemma~\ref{lem:monomial} and Proposition~\ref{prop:tower} and would approximate $\theta$ arbitrarily well. As $K$ is perfect, $\theta$ is separable; let $\kappa$ be the largest value of $v(\theta-\theta')$ over the other conjugates $\theta'$. A truncation $a\in K$ with $v(\theta-a)>\kappa$ would give $K(\theta)\subseteq K(a)=K$ by Krasner's lemma, a contradiction. A well-ordered subset of $\R$ cannot accumulate from the right, and its set of finite accumulation points is closed, so a least one $\delta$ exists; below every $r<\delta$ the support is finite and below $\delta$ it is infinite. Proposition~\ref{prop:tower} gives $\delta\in\Gamma$.
\end{proof}

\begin{lemma}\label{lem:cut}
For every $a\in K$ we have $v(\theta-a)<\delta$, and these values are cofinal below $\delta$.
\end{lemma}

\begin{proof}
If $v(\theta-a)\ge\delta$, then $a$ and $\theta$ agree below $\delta$ and $\supp a$ contains the infinite bounded set $\supp\alpha$, which is impossible for elements of $K$. For $r<\delta$, a finite truncation of $\theta$ at a bound between $r$ and $\delta$ lies in $K$ and approximates $\theta$ beyond $r$.
\end{proof}

\begin{proposition}\label{prop:additive}
The element $z=y^{-\delta}\alpha$ satisfies $P(z)=d$ for some $d\in K$ and some additive polynomial $P(X)=\sum_j\mu_jX^{p^j}\in k[X]$ with $\mu_0\ne0$.
\end{proposition}

\begin{proof}
The powers $\theta,\theta^p,\theta^{p^2},\dots$ lie in the $n$-dimensional $K$-space $K(\theta)$, so there is a nontrivial relation $\sum_ia_i\theta^{p^i}=0$ with $a_i\in K$; discard the zero coefficients. Let $q_i=v(a_i)$, let $\lambda_i\in k^\times$ be the leading coefficient of $a_i$, and put
\[
c=\min_i\bigl(q_i+p^i\delta\bigr),\qquad I=\{i:q_i+p^i\delta=c\}.
\]
Then $c\in\Gamma$. Write $\theta=\alpha+\eta$ with $v(\eta)\ge\delta$. Since Frobenius is additive, $a_i\theta^{p^i}=a_i\alpha^{p^i}+a_i\eta^{p^i}$ with $v(a_i\eta^{p^i})\ge c$, so the tails do not contribute below $c$.

Expand each $a_i$ into monomials $a_{i,q}y^q$. The term $a_{i,q}y^q\alpha^{p^i}$ has support in $(q+p^iv(\alpha),\,q+p^i\delta)$, so only finitely many $q$ contribute below $c$, because $a_i$ has finitely many support exponents below any bound. Since $q\ge q_i$ and $q_i+p^i\delta\ge c$, a contributing monomial has $q+p^i\delta\ge c$, with equality exactly when $i\in I$ and $q=q_i$. In that case the whole term $\lambda_iy^{q_i}\alpha^{p^i}$ lies below $c$. In every other case $q+p^i\delta>c$, and the part of $y^q\alpha^{p^i}$ below $c$ comes from a proper lower truncation of $\alpha$, which is finite. All these finite terms have exponents in $\Gamma$ and so lie in $K$ by Lemma~\ref{lem:monomial}. Truncating the relation below $c$ therefore gives
\[
\sum_{i\in I}\lambda_iy^{q_i}\alpha^{p^i}=R\in K .
\]
As $q_i=c-p^i\delta$, dividing by $y^c\in K$ gives $Q(z)=y^{-c}R\in K$ with $Q(X)=\sum_{i\in I}\lambda_iX^{p^i}\ne0$. If $m=\min I$, perfection of $k$ gives $P$ with nonzero linear coefficient and $Q=P^{p^m}$, perfection of $K$ gives $d\in K$ with $d^{p^m}=y^{-c}R$, and injectivity of Frobenius on $H$ gives $P(z)=d$.
\end{proof}

No algebraicity of $\alpha$ is assumed; the equation proves it. This extraction from an Ore relation follows the minimal-face argument of Kedlaya~\cite[Lemma~9.3]{Ked}; the specialization used here is that all truncated terms other than the leading ones are finite.

\begin{proposition}\label{prop:block}
The field $E=K(\alpha)=K(z)$ is a nontrivial elementary abelian Galois extension of $K$ contained in $K(\theta)$, and distinct $K$-conjugates of $\alpha$ differ by elements of valuation $\delta$.
\end{proposition}

\begin{proof}
The roots of $P$ form a finite $\F_p$-subspace $C$ of $k$, because $k$ is algebraically closed and $P$ is separable and additive, and the roots of $P(X)-d$ are $z+C\subseteq K(z)$. So the minimal polynomial of $z$ splits into distinct linear factors over $K(z)$ and $E/K$ is Galois. The map $\sigma\mapsto\sigma(z)-z$ from $\Gal(E/K)$ to $C$ is an injective homomorphism, since every $\sigma$ fixes $k$, so the group is elementary abelian. The extension is nontrivial because $\alpha$ has infinite bounded support, which no element of $K$ has. Two distinct conjugates of $\alpha$ differ by $y^\delta$ times a nonzero element of $k$.

For containment, let $b=\theta_\delta y^\delta\in K$, with $\theta_\delta=0$ if $\delta\notin\supp\theta$. Then $v(\theta-b-\alpha)>\delta$, while all conjugates of $\alpha$ lie at distance exactly $\delta$ from $\alpha$. Krasner's lemma gives $K(\alpha)\subseteq K(\theta-b)=K(\theta)$.
\end{proof}

\section{Proof of Theorem~\ref{thm:main}}\label{sec:proof}

Let $\theta$ be a primitive generator of $M/K$ of depth one, and let $\alpha,\delta,E$ be as above. If $[E:K]<[M:K]$, then $\alpha$ is algebraic of degree less than $\deg_K\theta$ and
\[
v(\theta-\alpha)\ge\delta>v(\theta-a)\qquad\text{for every }a\in K
\]
by Lemma~\ref{lem:cut}, contradicting~\eqref{eq:depthone}. Hence $[E:K]=[M:K]$, and since $E\subseteq K(\theta)=M$ we get $E=M$. By Proposition~\ref{prop:block}, $M/K$ is elementary abelian of exponent $p$, so it is the compositum of its cyclic subextensions of degree $p$. In characteristic $p$ each of these is an Artin--Schreier extension: if $\sigma$ generates its group, the trace of $1$ is $p=0$, additive Hilbert~90 gives $\gamma$ with $\sigma(\gamma)-\gamma=1$, and $\gamma^p-\gamma$ is fixed by $\sigma$. \qed

The argument uses only one primitive generator realizing depth one, not a particular generator of the tower. A second route, which only uses $\supp\theta\subseteq\Gamma$, obtains $\delta\in\Q$ from the extraction itself (equality of two of the affine functions $q_i+p^i\delta$), works over the tame extension $T=K(y^\delta)$, and descends to $K$ through the decomposition $\Gal(NT/K)\cong\Gal(N/K)\times\Gal(T/K)$ for a Galois $p$-extension $N$ containing $M$, using a normalized trace over $\Gal(T/K)$; we record it only as a check.

\section{Formal verification}\label{sec:lean}

Two Lean projects are included. The second states the theorem with depth as Nart and Novacoski define it; the first expresses depth one through Okutsu sequences.

\paragraph{The Okutsu formulation.} The project in \code{lean/} pins Lean~4.19.0 and the Mathlib commit
\begin{center}\code{c44e0c8ee63ca166450922a373c7409c5d26b00b}.\end{center}
Its 20 modules work with the actual Hahn field \code{HahnSeries} with rational exponents and coefficients in $k$. They construct Artin--Schreier roots by summable Hahn families, prove the support closure invariant of Proposition~\ref{prop:tower}, the first accumulation cut, the Ore extraction of Proposition~\ref{prop:additive}, the Galois and containment statements of Proposition~\ref{prop:block} with Krasner's lemma in an arbitrary ambient field, and the lower-degree obstruction of Section~\ref{sec:proof}.

The main endpoint is \code{okutsu_minimum_one_iff_asCompositum} in \code{HahnClosureForward.lean}. Its hypotheses are that $k$ is algebraically closed of characteristic $p$; that $K$ is an intermediate field of $k\subseteq H$ which is perfect, contains $y$ and consists of left-finite series; that the restricted valuation satisfies \code{RationalValuedHenselian}, which says that it has a unique extension with values in $\Q\cup\{\infty\}$ to every algebraic extension; and that $M\ne K$ is the top of a literal Artin--Schreier tower inside $H$. It works in the algebraic closure \code{hahnClosure K} of $K$ inside $H$, with the Hahn valuation, exactly as in~\cite[\S3]{NN}, and proves that $1$ is the least length of an Okutsu sequence over all primitive generators of $M$ if and only if the Artin--Schreier compositum of $M$ is all of $M$. The two implications are \code{asCompositum_of_hahnClosure_okutsu_minimum} (the converse, Theorem~\ref{thm:main}) and \code{hahnClosure_okutsu_minimum_of_asCompositum} (the forward implication, \cite[Theorem~3.15]{NN}, for any nontrivial finite-dimensional $M$ inside $H$ generated by Artin--Schreier roots).

The forward proof follows~\cite{NN}. Each Artin--Schreier generator is moved by an element of $K$ to the explicit negative root of a finite negative element of $K$; this uses henselian splitting of elements of nonnegative valuation, proved with Krasner's lemma. A combination of these roots with coefficients in $k$ that avoid finitely many linear relations is a primitive element whose nontrivial conjugates differ from it by nonzero constants. Krasner's lemma, after a constant shift in the boundary case, shows that no element of smaller degree comes within valuation zero of it, while partial inverse-Frobenius sums in $K$ approximate it at every negative level. An earlier endpoint, \code{asCompositum_of_algebraicClosure_okutsu_minimum}, proves the converse with Okutsu sequences in an abstract algebraic closure of $K$.

In this project the identification of the Okutsu minimum with depth is cited from~\cite[Theorem~1.3]{NN}, and henselianity enters as the unique extension of the valuation to algebraic extensions~\cite[\S4.1]{EP}.

\paragraph{The Mac Lane--Vaquié formulation.} The project in \code{lean\_mlv/} pins Lean~4.35.0-rc2 and the Mathlib commit \code{065356127b1d}. Its endpoint \code{artinSchreierTower_intrinsicMLVDepth_one_iff} in \code{ASDepthTheorem.lean} assumes that $k$ is algebraically closed of characteristic $p$, that $K$ is a perfect intermediate field of $k\subseteq H$ containing $y$ and consisting of left-finite series, and that the valuation ring of the restricted Hahn valuation is a henselian local ring in Mathlib's sense. For every finite tower $K=F_0\subseteq\cdots\subseteq F_N=E\ne K$ of abstract fields with $F_{i+1}=F_i(x_i)$ and $x_i^p-x_i\in F_i$, every $K$-embedding $\xi$ of $E$ into an algebraic closure of $K$ and every valuation $w$ on that closure extending the Hahn valuation of $K$, it proves that $1$ is the least length of a complete finite Mac Lane--Vaquié chain ending at evaluation at a primitive generator of $\xi(E)$ if and only if $E$ is generated over $K$ by its elements $z$ with $z^p-z\in K$. A companion endpoint constructs $\xi$ and $w$, and two further endpoints state the result with the family presentation of~\cite[Definitions~1.9--1.11]{N24}. The chains consist of key polynomials and ordinary and limit augmentations with their full expansion formulas, as in~\cite{N21}. Finite dimensionality, separability, the realization in $H$ and the existence of the extending valuation are derived, not assumed. The project vendors 54 modules of the NYLL library of N.~Yamaguchi for the existence and uniqueness of extensions of henselian valuations.

This project compares Mac Lane--Vaquié chains with the Okutsu criterion only at length one, which is all the theorem needs, and proves no equivalence at other lengths. The encodings of initial forms by strict cancellation and of the value group inside $\Q$ are explained in its documentation, not derived from a general library. In both projects the verifier \code{scripts/verify_cached.py} compiles every module into a fresh directory and checks that every audited declaration, including private and generated ones, depends only on \code{propext}, \code{Classical.choice} and \code{Quot.sound}. Lean verification establishes only the formal statements described here.

\section{Prior work and limits}\label{sec:novelty}

Nart and Novacoski introduced depth for defect extensions, proved that composita of Artin--Schreier extensions have depth one and stated Theorem~\ref{thm:main} as Conjecture~4.1 of~\cite{NN}, supporting it with a cyclic tower of depth two. The Artin--Schreier root formula, the support estimates, Krasner's lemma and the Ore relation are standard, and the minimal-face extraction is Kedlaya's~\cite[Lemma~9.3]{Ked}. What is offered is the first-block statement, Proposition~\ref{prop:block}, over these bases and its application to all towers and all primitive generators. Searches in October 2026 found no other proof of the conjecture, but a search cannot establish that none exists, and the published journal version of~\cite{NN} could not be opened. No priority is claimed, and no expert has reviewed the result.

\section*{Research process and AI disclosure}
This is fully autonomous agentic research. Dev D. Goyal is responsible for publication. No expert validation or external peer review is claimed, and the novelty assessment is provisional.

The theorem, its proof, the Lean formalization, the literature searches and an earlier version of this note were produced by autonomous AI research agents in an earlier session. The records of that session do not name the AI system used. A later session of autonomous agents, continuing the same project, produced the Lean~4.35 project: it formalized Mac Lane--Vaquié chains and their comparison at length one, replaced the unique-extension hypothesis by henselianity of the valuation ring and proved the result for abstract towers. Its records do not name the AI system used either. That project adapts Lean sources of the Tau Ceti project and of Y.~Yuan's TrustworthyKedlaya and vendors NYLL, whose author credits assistance by OpenAI Codex; the credits are in \code{lean\_mlv/THIRD\_PARTY.md}.

Opus 5.5 by Anthropic used with Claude Code assessed the package for this release. It checked the argument step by step, read the arXiv source of~\cite{NN} (version~2) and checked the depth criterion~\eqref{eq:depthone}, Theorem~1.3, Theorem~3.15 and Conjecture~4.1 against the statements used here, and compared the Lean endpoints and their hypotheses with Theorem~\ref{thm:main}. It formalized the converse in the source's ambient field, the forward implication and their equivalence (the modules \code{HahnClosureDepth} and \code{HahnClosureForward}), ran the fresh Lean replay described in Section~\ref{sec:lean}, and rewrote the note in this form. For the Lean~4.35 project it compared the endpoints and the chain definitions with Theorem~\ref{thm:main} and~\cite{N21}, ran its fresh replay and revised this note. AI checks are not human peer review. Lean verification establishes only the formal statements described in Section~\ref{sec:lean}.

\begin{thebibliography}{9}\small
\bibitem{EP} A.~J.~Engler and A.~Prestel, \emph{Valued Fields}, Springer, 2005.
\bibitem{Ked} K.~S.~Kedlaya, On the algebraicity of generalized power series, \emph{Beitr.\ Algebra Geom.} \textbf{58} (2017), 499--527. \href{https://arxiv.org/abs/1508.01836}{arXiv:1508.01836}.
\bibitem{N21} E.~Nart, MacLane--Vaquié chains of valuations on a polynomial ring, \emph{Pacific J. Math.} \textbf{311} (2021), 165--195. \href{https://arxiv.org/abs/1911.01714}{arXiv:1911.01714}.
\bibitem{N24} E.~Nart, Okutsu sequences in henselian valued fields, \emph{Polynesian J. Math.} \textbf{1} (2024), no.~4.
\bibitem{NN} E.~Nart and J.~Novacoski, Depth of Artin--Schreier defect towers, \emph{J. Pure Appl.\ Algebra} \textbf{230} (2026), 108184. \href{https://doi.org/10.1016/j.jpaa.2026.108184}{doi:10.1016/j.jpaa.2026.108184}; \href{https://arxiv.org/abs/2503.18827}{arXiv:2503.18827v2} (16 October 2026).
\end{thebibliography}

\end{document}
